\subsection{阶函数}

	设 A 为滤过环，E 为滤过 A-模，\((\mathrm{E}_{n})\) 为 E 的滤过。对所有 \(x \in E\)，令 \(v(x)\) 表示这些整数在 \(\bar{R}\) 中的上确界，其中这些整数为 \(n \in Z\)，并满足 \(x \in E_{n}\)。则有下列等价关系：

\[
\left\{ \begin{array}{l} v (x) = - \infty \Leftrightarrow x \notin \bigcup_ {n \in \mathbf {Z}} \mathrm{E} _ {n} \\ v (x) = p \quad \Leftrightarrow x \in \mathrm{E} _ {p} \quad \text { and } \quad x \notin \mathrm{E} _ {p + 1} \\ v (x) = + \infty \Leftrightarrow x \in \bigcap_ {n \in \mathbf {Z}} \mathrm{E} _ {n} \end{array} \right.\tag{4}
\]

	映射 \(v: \mathrm{E} \to \overline{\mathbf{R}}\) 称为滤过模 \(\mathrm{E}\) 的阶函数。如果已知 \(v\)，则各个 \(\mathrm{E}_n\) 也随之确定，因为 \(\mathrm{E}_n\) 是由满足 \(x \in \mathrm{E}\) 且 \(v(x) \geqslant n\) 的元素构成的集合；\(\mathrm{E}_n\) 是 \(\mathrm{E}\) 的加法子群这一事实给出关系式

\[
v (x - y) \geqslant \inf (v (x), v (y)).\tag{5}
\]

上述定义尤其适用于滤过 A-模 \(A_{s}\)；令 w 为其阶函数。由第 1 节的公式 (3) 可知，对 \(a \in A\) 和 \(x \in E\)，

\[
v (a x) \geqslant w (a) + v (x)\tag{6}
\]

只要右端有定义；特别地，对 \(a \in A\) 和 \(b \in A\)，

\[
w (a b) \geqslant w (a) + w (b)\tag{7}
\]

只要右端有定义。

对于不必交换的滤过群 G，也类似地定义阶函数；这时与 (5) 相应的关系写成

\[
v (y x ^ {- 1}) = v (x y ^ {- 1}) \geqslant \inf (v (x), v (y)).\tag{5'}
\]

